\documentclass[12pt]{article}
\usepackage{amsmath,amssymb,amsfonts,amsthm}
\usepackage{braket}
\usepackage{hyperref}
\usepackage{geometry}
\usepackage{authblk}
\geometry{margin=1in}

\newtheorem{theorem}{Theorem}[section]
\newtheorem{lemma}[theorem]{Lemma}
\newtheorem{proposition}[theorem]{Proposition}
\newtheorem{corollary}[theorem]{Corollary}
\theoremstyle{definition}
\newtheorem{definition}[theorem]{Definition}
\newtheorem{assumption}[theorem]{Assumption}
\theoremstyle{remark}
\newtheorem{remark}[theorem]{Remark}

\newcommand{\Hthree}{\mathbb{H}^3}
\newcommand{\Sphere}{\mathbb{S}^2}
\newcommand{\GammaG}{\Gamma}
\newcommand{\U}{\mathrm{U}}
\newcommand{\calH}{\mathcal{H}}
\newcommand{\bdry}{\partial\mathbb{H}^3}
\newcommand{\calHS}{\mathcal{H}_S}
\newcommand{\calHT}{\mathcal{H}_T}

\title{Flavor Mixing from Boundary Misalignment\\in Hyperbolic Geometry}
uthor[1]{Marvin L.\ Gentry}
ffil[1]{Independent Researcher\
  	exttt{drmlgentry@protonmail.com}\
  ORCID: 0009-0006-4550-2663}
\date{}

\begin{document}
\maketitle

\begin{abstract}
We formalize a kinematic mechanism by which fermion mixing matrices arise
from the misalignment between discrete boundary sectors associated with
logarithmic mass organization in hyperbolic space. Mixing matrices are
defined as overlap operators between preferred orthonormal bases of sector
subspaces and are shown to be uniquely determined up to sector-local
rephasings. Our main result (Theorem~\ref{thm:discrete-classes}) establishes
that the set of accessible mixing matrices is discrete in $\U(n)$ modulo
sector-local rephasings, providing a structural explanation for the
stability of mixing patterns without continuous parameter tuning.
\end{abstract}

\section{Introduction}

The Standard Model fermion mixing matrices exhibit stable, non-generic
patterns unexplained by gauge structure alone. We develop a kinematic
explanation: mixing arises from the misalignment between boundary sectors
naturally associated with discrete logarithmic mass organization.

If fermion mass scales are organized near a discrete logarithmic lattice,
each fermion is associated with a definite direction on the boundary sphere
$\bdry\cong\Sphere$. Different generations of a given sector span different
subspaces of a boundary Hilbert space $\calH$; their overlap is the mixing
matrix. The key structural result is that if boundary directions are assigned
by a \emph{discrete} selection rule, the set of accessible mixing matrices
is itself discrete (Theorem~\ref{thm:discrete-classes}).

\section{Geometric Setup and the Discreteness Assumption}

\begin{definition}[Boundary direction]\label{def:bdrydir}
A \emph{boundary direction} is a point $\xi\in\bdry\cong\Sphere$.
Group elements $\gamma\in\GammaG$ are \emph{associated with} $\xi$ when
the geodesic ray from basepoint $0$ through $\gamma\cdot0$ converges to~$\xi$.
\end{definition}

\begin{assumption}[Discrete sector assignment]\label{ass:discrete}
Each fermion flavor $f\in F$ is assigned a unique boundary direction
$\xi(f)\in\bdry$ by a deterministic, discrete selection rule: the image
$\Xi=\{\xi(f):f\in F\}$ is finite, and the assignment is locally rigid.
One concrete realization: $\xi(f)$ is the ideal endpoint of the geodesic
ray to the minimal representative $\gamma_f\cdot0$, whose uniqueness andrigidity follow from the companion paper on operator realization~\cite{bridsonhaefliger}.
\end{assumption}

\section{Boundary Hilbert Space and Localized States}

\begin{definition}[Boundary-localized state]\label{def:localized}
For each $\xi\in\Sphere$, fix a normalized state $\ket{\xi}\in\calH$
with $\langle\xi|\xi\rangle=1$, depending continuously on $\xi$, with
$\langle\xi|\zeta\rangle=1$ iff $\xi=\zeta$.
\end{definition}

\section{Sector Subspaces and Preferred Bases}

\begin{definition}[Sector subspace]\label{def:sectorspan}
The \emph{sector subspace} of $S$ is
$\calHS=\overline{\mathrm{span}}\{\ket{\xi(f)}:f\in S\}\subset\calH$,
with $\dim\calHS=|S|$ by linear independence.
\end{definition}

\begin{definition}[Preferred orthonormal basis]\label{def:preferred}
The \emph{preferred ONB} of $\calHS$ is the Gram--Schmidt
orthonormalization of $(\ket{\xi(f_1)},\ldots,\ket{\xi(f_{n_S})})$
in the fixed sector ordering, with all pivot coefficients real and positive.
\end{definition}

\section{Mixing Matrices as Overlap Operators}

\begin{definition}[Mixing matrix]\label{def:mixing}
For sectors $S$, $T$ of equal cardinality $n$ with preferred ONBs
$\{\ket{S,i}\}$ and $\{\ket{T,j}\}$, the \emph{mixing matrix} is
$(U_{ST})_{ji}=\braket{T,j|S,i}$.
\end{definition}

\begin{proposition}[Unitarity]\label{prop:unitary}
If $\calHS=\calHT$, then $U_{ST}\in\U(n)$.
\end{proposition}

\begin{theorem}[Geometric dependence]\label{thm:geom}
$U_{ST}$ is determined entirely by $\Xi_{ST}=\{\xi(f):f\in S\cup T\}$
and the inner product on $\calH$.
\end{theorem}

\section{Rephasing Invariance}

\begin{proposition}[Rephasing-invariant quantities]\label{prop:rephinv}
Under sector-local rephasings $\ket{S,i}\mapsto e^{i\alpha_i}\ket{S,i}$,
$\ket{T,j}\mapsto e^{i\beta_j}\ket{T,j}$, the following are invariant:
\emph{(i)} the moduli $|(U_{ST})_{ji}|^2$;
\emph{(ii)} the quartic products
$J_{ji,j'i'}=(U_{ST})_{ji}(U_{ST})_{j'i'}\overline{(U_{ST})_{ji'}}\,\overline{(U_{ST})_{j'i}}$;
\emph{(iii)} $\det(U_{ST})$ up to overall phase.
\end{proposition}

\section{Discreteness and Stability of Mixing Classes}

\begin{theorem}[Stability]\label{thm:stability}
For $\varepsilon\leq\varepsilon_0=\delta_{\min}/2$, where
$\delta_{\min}=\min_{f\neq f'}d_\Sphere(\xi(f),\xi(f'))$, any bounded
deformation of magnitude $\varepsilon$ leaves $\Xi$, the preferred bases,
and $U_{ST}$ unchanged.
\end{theorem}

\begin{theorem}[Discreteness of mixing classes]\label{thm:discrete-classes}
Under Assumption~\ref{ass:discrete}, the set of accessible mixing matrices
$\mathcal{M}_{ST}$ is discrete in $\U(n)/(\U(1)^n\times\U(1)^n)$.
\end{theorem}

\begin{proof}
The set of admissible configurations $\Xi_{ST}$ is finite (discrete).
By Theorem~\ref{thm:geom}, the map $\Xi_{ST}\mapsto U_{ST}$ is continuous,
so it maps a finite set to a finite set. Quotienting by rephasings
preserves finiteness.
\end{proof}

\begin{corollary}[No continuous drift]\label{cor:nodrift}
No continuous deformation of physical parameters can produce continuous
drift in any rephasing-invariant mixing observable.
\end{corollary}

\section{Conclusion}

We have shown that fermion mixing matrices arising from boundary-sector
misalignment are unitary by construction, rephasing-invariant, stable under
bounded deformations, and discrete modulo rephasings. These properties are
structural consequences of the discrete selection rule and require no
dynamical input or continuous tuning.

\bibliographystyle{plain}
\bibliography{gentry-flavor-mixing}

\end{document}